% !TeX program = xelatex
\PassOptionsToPackage{quiet}{fontspec}
\documentclass[12pt,a4paper,UTF8]{ctexart}
\CTEXsetup[format={\Large\bfseries}]{section}
\usepackage{zjureport}
\usepackage{fontawesome}

%--------------------信息设置---------------------------%
\mytitle{短时傅里叶变换可视化}
\myauthor{xx}
\mystudentid{xx}
\mycourse{信号与系统}
\myschool{信息与电子工程学院}

\begin{document}
\cover
\thispagestyle{empty}

\newpage
\setcounter{page}{1}

\begin{abstract}
傅里叶变换能够揭示信号的全局频率组成，但无法回答"某个频率在什么时候出现"这一时变分析问题。针对该局限，本文围绕短时傅里叶变换（Short-Time Fourier Transform, STFT）展开课程设计。首先，通过 \texttt{waveform.py} 绘制两路离散信号的时域波形，再利用 \texttt{fft\_analyze.py} 调用手写数值傅里叶变换验证两者在全局频谱上几乎一致；随后基于 \texttt{stft\_analyzer.py} 实现的滑动矩形窗逐帧变换，绘制不同窗长下的时频图，分析时间分辨率与频率分辨率之间的权衡；最后以自选音乐为对象，利用 \texttt{stft\_animator.py} 与 \texttt{fft\_animator.py} 分别生成滚动时频动画与彩虹雨频谱动画，并将所有分析功能封装为 Flask 后端 API，结合 Vue.js 前端部署为可交互网站。实验结果表明：仅观察频谱图无法区分局部时序结构，而 STFT 可有效揭示两路信号在频率出现顺序上的差异；窗长减小时时频图边界更清晰但频率纹理变粗；对于音乐信号，STFT 能直观反映主旋律、谐波和瞬态打击成分的时间演化。
\end{abstract}

\noindent\textbf{关键词：}短时傅里叶变换；时频分析；矩形窗；音频可视化；Web 部署

\section{问题提出与研究意义}
经典傅里叶变换通过对整个时间轴进行积分，将时域信号映射到频域，从而揭示信号所包含的频率分量。然而这种变换是全局性的：一旦频率成分随时间发生变化，普通频谱图只能给出"包含了哪些频率"，却无法说明"这些频率分别在什么时刻出现、持续了多久"。对于语音、音乐以及非平稳工程信号，这一缺陷会显著限制分析能力。

本课程设计围绕"短时傅里叶变换可视化"这一主题展开，目标包括四个方面：
\begin{enumerate}
    \item 对给定的两路离散信号进行时域观察与傅里叶变换，比较其全局频谱特性；
    \item 采用矩形窗自行实现 STFT，绘制两路信号在不同窗长下的时频图，讨论时间分辨率与频率分辨率之间的权衡；
    \item 选择一段音乐，分别生成滚动时频动画（STFT 动画）和逐帧频谱动画（FFT 动画），实现音频的动态可视化；
    \item 将所有分析功能封装为 Web 应用，在线部署，使用户能够上传任意音频文件并实时获取时频分析结果。
\end{enumerate}

该课题将"全局频谱分析"与"局部时频分析"联系起来，有助于理解非平稳信号处理中最核心的时频分辨率权衡问题，也为后续学习语音处理、音乐信息检索和雷达声纳信号分析提供基础。

\section{实验原理与实现方法}
\subsection{连续傅里叶变换与离散实现}
连续时间信号 $x(t)$ 的傅里叶变换定义为
\[
X(f)=\int_{-\infty}^{\infty} x(t)e^{-j2\pi ft}\,\mathrm{d}t .
\]
题目给出的信号已按采样率 $f_s=4410\ \mathrm{Hz}$ 离散化，将积分改写为数值求和。若信号长度为 $N$，采样时刻为 $t_n=n/f_s$，频率采样点为 $f_k=k f_s/N$，则：
\[
X(f_k)\approx \sum_{n=0}^{N-1}x[n]e^{-j2\pi f_k t_n}\cdot \frac{1}{f_s}.
\]
\texttt{py\_tools/lib/manual\_fft.py} 按上式实现了三个函数：\texttt{build\_time\_samples} 生成时间轴 $t_n$，\texttt{build\_frequency\_samples} 生成频率轴 $f_k$，\texttt{fft} 通过双重循环完成数值求和。\texttt{fft\_analyze.py} 调用上述函数，从 \texttt{.mat} 文件中读取信号变量（优先取 \texttt{z}，兼容 \texttt{zz}）并将幅值谱保存为图像；\texttt{waveform.py} 则负责绘制时域波形，两者的输出均保存至 \texttt{py\_tools/output}。

\subsection{短时傅里叶变换}
为了获得频率随时间的演化规律，需要用有限长度的窗函数 $w(t)$ 在时间轴上滑动，每移动到一个位置就截取局部片段并进行傅里叶变换。连续形式为
\[
\mathrm{STFT}_x(t,f)=\int_{-\infty}^{\infty}x(\tau)w(\tau-t)e^{-j2\pi f\tau}\,\mathrm{d}\tau .
\]
离散实现可写成
\[
X_m[k]=\sum_{n=0}^{L-1}x[n+mR]\,w[n]\,e^{-j2\pi kn/L},
\]
其中 $L$ 为窗长，$R$ 为帧移，$m$ 为帧序号。题目要求第二部分使用矩形窗：
\[
w[n]=
\begin{cases}
1, & 0\le n < L,\\
0, & \text{其他}.
\end{cases}
\]

STFT 的核心约束在于时间分辨率与频率分辨率不能同时任意提高。窗长为 $L$ 时，频率分辨率近似为
\[
\Delta f=\frac{f_s}{L},
\]
帧移为 $R$ 时，时间分辨率近似为
\[
\Delta t=\frac{R}{f_s}.
\]
较小的窗长能够更清楚地定位频率变化发生的时刻，但频率刻画变粗；较大的窗长提高频率分辨率，但会模糊时间上的突变边界。

\texttt{stft\_analyzer.py} 中的 \texttt{stft} 函数对每一帧直接调用 \texttt{numpy.fft.fft}，而非复用手写的双重循环。这是因为 STFT 需要对大量重叠帧重复计算频谱，FFT 将单帧计算复杂度降为 $O(L\log L)$，远低于双重循环的 $O(L^2)$，更适合批量时频图绘制和长音频动画生成。

\subsection{代码组织与实验流程}
项目代码按功能分为两个部分：\texttt{py\_tools} 目录存放离线分析脚本，\texttt{signal\_analyzer} 目录存放 Web 部署代码。

\textbf{离线分析脚本（\texttt{py\_tools}）}共包含以下模块：
\begin{enumerate}
    \item \texttt{lib/manual\_fft.py}：提供手写数值傅里叶变换函数 \texttt{fft}，以及配套的时间轴生成函数 \texttt{build\_time\_samples} 和频率轴生成函数 \texttt{build\_frequency\_samples}；
    \item \texttt{signal\_stats.py}：计算两路信号的均值、方差和总能量，验证它们在全局统计特性上的一致性；
    \item \texttt{waveform.py}：读取 \texttt{signal1.mat} 和 \texttt{signal2.mat}，通过 \texttt{Config} 类统一管理图像尺寸（$12\times4$ 英寸）和输出路径，绘制并保存时域波形图；
    \item \texttt{music\_waveform.py}：分析音乐文件的波形图；
    \item \texttt{fft\_analyze.py}：调用 \texttt{manual\_fft} 计算两路信号的全局频谱，读取 \texttt{.mat} 变量时兼容 \texttt{z} 与 \texttt{zz} 两种命名，将幅值谱图保存到 \texttt{py\_tools/output}；
    \item \texttt{music\_fft.py}：分析音乐信号频谱图
    \item \texttt{stft\_analyzer.py}：实现带矩形窗的 STFT 核心函数 \texttt{stft}，默认帧移为半窗重叠（$R=L/2$），批量输出不同窗长下的时频图；
    \item \texttt{music\_stft.py}：分析音乐信号的时频图
    \item \texttt{stft\_animator.py}：读取 \texttt{self\_songbie.mp3}，复用 \texttt{stft\_analyzer.py} 的 \texttt{stft} 函数生成滚动时频动画，每帧展示宽度为 10 帧的滑动时间窗口，使时频图随音乐连续滚动，由 \texttt{FFMpegWriter} 导出静音 MP4 后再调用 \texttt{ffmpeg} 混入原始音轨；
    \item \texttt{fft\_animator.py}：逐帧对音频信号做矩形窗 FFT，将 dB 标度频谱绘制为折线（对数频率轴），同时在背景随机生成 120 条 HSV 彩色线段叠加为"彩虹雨"装饰效果，生成视觉更丰富的频谱动画；
    \item \texttt{lib/log.py}：轻量日志模块，同时输出至控制台和 \texttt{py\_tools/data/project.log}；
    \item \texttt{py\_tools/output}：统一保存所有频谱图、时频图和视频输出，便于后续插图整理与实验对比。
\end{enumerate}

\textbf{Web 应用（\texttt{signal\_analyzer}）}采用前后端分离架构：后端由 Python Flask 提供 REST API，前端由 Vue.js 实现交互界面，通过 Docker Compose 统一部署。后端核心模块如下：
\begin{enumerate}
    \item \texttt{services/stft\_analyzer.py}：STFT 引擎，在离线版矩形窗基础上增加了 Hann、Hamming、Blackman、Bartlett、Kaiser 五种窗函数，并支持独立设置 FFT 点数（零填充），接口签名与离线版保持一致；
    \item \texttt{services/signal\_processing.py}：接收 \texttt{AnalysisParams} 字典（含 \texttt{sr}、\texttt{n\_fft}、\texttt{hop\_length}、\texttt{window} 等字段），返回波形图、频谱图和时频图的 Base64 编码字符串；
    \item \texttt{services/stft\_animation.py}：接收 \texttt{AnimationParams} 字典，生成滚动时频动画 MP4；
    \item \texttt{api/signal.py}：定义 \texttt{POST /upload}、\texttt{GET /analyze/<filename>}、\texttt{POST /animate/<filename>} 等 REST 接口，从查询字符串或请求体中读取可调参数，动画渲染通过异步任务队列完成，避免长时间阻塞请求。
\end{enumerate}

materials文件夹内部存放分析的数据和分析出的结果。

signal1.mat  原始信号1

signal2.mat  原始信号2

self\_songbie.mp3  时频动画的原始音乐文件

music\_stft.mp4  音乐的滚动时频动画

music\_fft.mp4  音乐的彩虹雨频谱动画

据此，全文实验流程可概括为三个阶段：首先，运行 \texttt{waveform.py} 和 \texttt{fft\_analyze.py} 获得两路给定信号的时域波形与全局频谱，建立"频率成分基本一致"的基线结论；其次，运行 \texttt{stft\_analyzer.py}，在多组窗长条件下比较两路信号的时频图，分析分辨率折中关系；最后，运行 \texttt{stft\_animator.py} 和 \texttt{fft\_animator.py} 对音乐片段生成两种形式的动画，并通过 \texttt{signal\_analyzer} 将全部功能部署为线上交互服务。前两组实验的输入均为长度约 $14000$ 点的 \texttt{.mat} 信号，采样率 $f_s=4410\ \mathrm{Hz}$；音乐实验的输入为 \texttt{self\_songbie.mp3}，按 $44.1\ \mathrm{kHz}$ 读取处理。

\section{给定信号的频谱特性分析}
图 \ref{fig:time-domain} 给出了两路信号的时域波形。可以看到，二者的局部振荡方式和能量分布明显不同，仅从波形外观就可以判断它们并非同一个信号。

\begin{figure}[H]
    \centering
    \subfigure[信号1时域波形]{
        \includegraphics[width=0.48\textwidth]{waveforms/signal1_time.png}
    }
    \hfill
    \subfigure[信号2时域波形]{
        \includegraphics[width=0.48\textwidth]{waveforms/signal2_time.png}
    }
    \caption{给定信号的时域波形对比}
    \label{fig:time-domain}
\end{figure}

随后利用 \texttt{py\_tools/fft\_analyze.py} 计算两路信号的频谱，结果如图 \ref{fig:spectrum} 所示。

\begin{figure}[H]
    \centering
    \subfigure[信号1频谱图]{
        \includegraphics[width=0.48\textwidth]{spectrum/signal1.png}
    }
    \hfill
    \subfigure[信号2频谱图]{
        \includegraphics[width=0.48\textwidth]{spectrum/signal2.png}
    }
    \caption{给定信号的傅里叶变换结果}
    \label{fig:spectrum}
\end{figure}

从图 \ref{fig:spectrum} 可以得到结论：

两路信号的主峰位置与幅值几乎完全重合，典型峰值出现在0-100Hz和4k-5kHz附近，两个峰值关于$\frac{f_s}{2}$对称分布。

进一步计算(py\_tools/signal\_stats.py)可知，两路信号的均值、方差和总能量也基本一致。这说明它们包含了几乎相同的全局频率成分，但这些频率成分在时间轴上的组织方式不同。仅凭普通傅里叶变换无法解释这种"频谱相同、波形不同"的现象，因此必须引入时频分析。
\begin{table}[H]
\centering
\caption{信号基本特征分析}
\begin{tabular}{|c|c|c|c|}
\hline
\textbf{Signal} & \textbf{Mean} & \textbf{Variance} & \textbf{Total Energy} \\
\hline
signal1.mat & 0.0049 & 1.0001 & 14001.6744 \\
\hline
signal2.mat & 0.0049 & 1.0001 & 14001.6744 \\
\hline
\end{tabular}
\end{table}
\section{短时傅里叶变换结果与窗长规律}
\subsection{实验设置}
本节展示的 STFT 结果均采用矩形窗（见 \texttt{py\_tools/stft\_analyzer.py}），并令帧移为半个窗长 $frame shift=window size/2$。为便于观察分辨率变化规律，正文选取 $window size=700$、$1400$、$2800$ 多种代表性窗长进行比较，相应参数如表 \ref{tab:stft-params} 所示。

\begin{table}[H]
    \centering
    \begin{tabular}{|c|c|c|c|}
        \hline
        窗长 $window size$ & 帧移 $frame shift$ & 频率分辨率 $\Delta f$ & 时间分辨率 $\Delta t$ \\
        \hline
        700 & 350 & $6.30\ \mathrm{Hz}$ & $79.4\ \mathrm{ms}$ \\
        \hline
        1400 & 700 & $3.15\ \mathrm{Hz}$ & $158.7\ \mathrm{ms}$ \\
        \hline
        2800 & 1400 & $1.58\ \mathrm{Hz}$ & $317.5\ \mathrm{ms}$ \\
        \hline
    \end{tabular}
    \caption{不同窗长下的 STFT 参数（$f_s=4410\ \mathrm{Hz}$）}
    \label{tab:stft-params}
\end{table}

\subsection{信号1的时频图}
下图 给出了信号1在三种窗长下的 STFT 结果。可以看到，信号1的主频随时间整体呈分段上升趋势：较短窗长时，频率变化发生的位置更容易被定位，时频图中的边界更清晰；较长窗长时，频率条纹更细、更稳定，但相邻时段之间的切换变得平滑。

\begin{figure}[H]
    \centering
    \includegraphics[width=0.4\textwidth, height=0.2\textheight]{stft/signal1_stft_W700.png}
    \caption{信号1在窗长700下的时频图}
\end{figure}
\begin{figure}[H]
    \centering
    \includegraphics[width=0.4\textwidth, height=0.2\textheight]{stft/signal1_stft_W1400.png}
    \caption{信号1在窗长1400下的时频图}
\end{figure}
\begin{figure}[H]
    \centering
    \includegraphics[width=0.4\textwidth, height=0.2\textheight]{stft/signal1_stft_W2800.png}
    \caption{信号1在窗长2800下的时频图}
\end{figure}

\subsection{信号2的时频图}
下图 给出了信号2在相同参数下的 STFT 结果。与信号1相比，信号2仍然由相近的几个主频成分构成，但这些频率成分在时间上的排列顺序已经改变，局部能量分布也明显不同。

\begin{figure}[H]
    \centering
    \includegraphics[width=0.4\textwidth, height=0.2\textheight]{stft/signal2_stft_W700.png}
    \caption{信号2在窗长700下的时频图}
\end{figure}
\begin{figure}[H]
    \centering
    \includegraphics[width=0.4\textwidth, height=0.2\textheight]{stft/signal2_stft_W1400.png}
    \caption{信号2在窗长1400下的时频图}
\end{figure}
\begin{figure}[H]
    \centering
    \includegraphics[width=0.4\textwidth, height=0.2\textheight]{stft/signal2_stft_W2800.png}
    \caption{信号2在窗长2800下的时频图}
\end{figure}
对比两组信号的时频图可以发现：
\begin{enumerate}
    \item 两路信号的全局频率集合基本相同，这与频谱图的结论一致；
    \item 信号1和信号2的主要频率成分的频率是逐步下降的，但是时间先后相互错位，对比700窗长的时频图非常明显。
    \item 当窗长从 700 增大到 2800 时，频率条纹变得更细（频率分辨率提升），同时时间边界逐渐模糊（时间分辨率下降），与表 \ref{tab:stft-params} 中的理论预测吻合。
\end{enumerate}

因此，STFT 成功揭示了普通频谱图无法体现的"局部时序差异"，说明对于非平稳信号，仅有全局频谱远远不够，必须联合时间轴观察频率的出现过程。

\section{音乐信号可视化分析}
\subsection{音乐信号处理流程}
本节以 \texttt{self\_songbie.mp3} 为输入，分别生成波形图、频谱图、时频图、滚动时频动画和彩虹雨频谱动画。三个脚本共享相同的参数配置：采样率 $f_s=44100\ \mathrm{Hz}$，FFT 点数 $N_\mathrm{fft}=2^{15}=32768$，帧移 $H=2^{10}=1024$。对应的分辨率为
\[
\Delta f=\frac{f_s}{N_\mathrm{fft}}=\frac{44100}{32768}\approx1.35\ \mathrm{Hz},\qquad
\Delta t=\frac{H}{f_s}=\frac{1024}{44100}\approx23.2\ \mathrm{ms}.
\]
此参数组合以牺牲极少的时间分辨率为代价，换取了极高的频率分辨率，适合展示音乐信号中细密的谐波结构。

三个脚本的具体处理流程如下：
\begin{enumerate}

    \item \textbf{波形图（\texttt{music\_stft.py}）}：使用 \texttt{librosa.load} 读取音频并单声道化，调用 \texttt{stft\_analyzer.py} 中的 \texttt{stft} 函数得到复数时频矩阵，保留正频率部分计算幅值并转换为 dB 标度，在对数频率轴下绘制时频图与时域波形的组合快照；
    \item \textbf{频谱图（\texttt{music\_fft.py}）}：对整个音频信号做一次矩形窗 FFT，计算幅值谱并转换为 dB 标度，在对数频率轴下绘制频谱图；
    \item \textbf{静态时频图（\texttt{music\_stft.py}）}：使用 \texttt{librosa.load} 读取音频并单声道化，调用 \texttt{stft\_analyzer.py} 中的 \texttt{stft} 函数得到复数时频矩阵，保留正频率部分计算幅值并转换为 dB 标度，在对数频率轴下绘制时频图与时域波形的组合快照；
    \item \textbf{时频动画（\texttt{stft\_animator.py}）}：在静态时频图基础上引入动态显示逻辑，每帧展示宽度为 10 帧的滑动时间窗口（约 $10\times23.2\ \mathrm{ms}=232\ \mathrm{ms}$），使时频矩阵的局部切片随音乐播放连续滚动；时间戳以文字形式叠加在画面右上角；由 \texttt{FFMpegWriter} 先导出静音 MP4，再调用 \texttt{ffmpeg} 将原始音轨混流进最终视频；
    \item \textbf{频谱动画（\texttt{fft\_animator.py}）}：每帧对当前音频片段做矩形窗 FFT，将 dB 标度的正频率频谱绘制为折线（对数横轴），同时在背景中随机生成 120 条 HSV 彩色线段作为"彩虹雨"装饰，线段的横坐标在对数频率轴上均匀分布，纵坐标随机落在 dB 范围内，帧率与 STFT 动画一致，最终同样导出带音轨的 MP4。
\end{enumerate}

\subsection{静态分析}

\begin{figure}[H]
    \centering
    \includegraphics[width=\textwidth]{waveforms/music_waveform.png}
    \caption{音乐信号的波形图}
    \label{fig:songbie waveform}
\end{figure}

\begin{figure}[H]
    \centering
    \includegraphics[width=\textwidth]{spectrum/music_fft.png}
    \caption{音乐信号的频谱图}
    \label{fig:songbie spectrum}
\end{figure}
\begin{figure}[H]
    \centering
    \includegraphics[width=\textwidth]{stft/music_stft_W32768.png}
    \caption{音乐信号的时频图}
    \label{fig:songbie stft}
\end{figure}

从上图可以观察到以下特点：
\begin{enumerate}
    \item 音乐信号的能量主要集中在中低频区域，主旋律与基频结构清晰可见；
    \item 时频图中出现多组稳定的水平谐波纹理，反映了乐音的音高结构；
    \item 在节拍或强起音附近，出现更宽的竖向能量扩展，对应打击感更强的瞬态成分。
\end{enumerate}

\subsection{时频动画分析}
时频动画（见 \texttt{py\_tools/stft\_animator.py}）通过 10 帧滑动窗口实时展示时频矩阵的局部切片，使观察者能够随着音乐推进直观感受频率能量的时间演化。与静态时频图相比，动画形式能够将音高的起伏走向、和声叠加时频率条纹的聚合与分散、以及节拍处亮度的骤变展现得更加直观。由于每帧只显示约 $232\ \mathrm{ms}$ 的时间窗口，局部频率结构的细节得以被放大，有助于分辨相邻音符之间的过渡。
\begin{figure}[H]
    \centering
    \includegraphics[width=0.8\textwidth]{stft_animation.png}
    \caption{时频动画截图}
    \label{fig:stft-animator}
\end{figure}
完整视频见附件materials/music\_stft.mp4
\subsection{频谱动画分析}
频谱动画（见 \texttt{py\_tools/fft\_animator.py}）以折线形式展示每一帧的瞬时 FFT 频谱，彩虹雨装饰背景使视觉效果更具感染力。由于每帧仅显示单帧 FFT 结果，与时频动画相比，频谱动画更直观地呈现了某一时刻各频率分量的相对强度；但缺少时间维度的二维展示，低频段谱线的起伏能够直接反映旋律的进行，高频段则可见伴奏和打击乐器的细节。
\begin{figure}[H]
    \centering
    \includegraphics[width=0.8\textwidth]{fft_animation.png}
    \caption{频谱动画截图}
    \label{fig:fft-animator}
\end{figure}
完整视频见附件materials/music\_fft.mp4
\subsection{Web 应用部署}
为了方便他人复现和使用上述分析功能，本项目基于 \texttt{signal\_analyzer} 模块构建了可交互网站，部署于魔搭社区（ModelScope）创空间，网址为：

\begin{center}
\texttt{https://www.modelscope.cn/studios/XueZhe4/signal-analyzer}
\end{center}
\begin{figure}[H]
    \centering
    \includegraphics[width=0.8\textwidth]{web.png}
    \caption{Web 应用界面截图}
    \label{fig:web}
\end{figure}
网站采用前后端分离架构，技术栈如下：
\begin{itemize}
    \item \textbf{后端}：Python Flask，提供 REST API；核心接口包括 \texttt{POST /upload}（文件上传）、\texttt{GET /analyze/<filename>}（同步返回波形图、频谱图、时频图的 Base64 编码）和 \texttt{POST /animate/<filename>}（异步提交动画渲染任务，返回任务 ID）；
    \item \textbf{前端}：Vue 3 + TypeScript + Vite，提供参数输入表单（采样率、FFT 点数、帧移、窗函数类型、颜色方案等可调参数）、图像展示面板和视频播放器；
    \item \textbf{部署}：Docker Compose 编排前后端容器，Nginx 反向代理前端静态资源，FFmpeg 集成于 Docker 镜像以支持视频渲染。
\end{itemize}

与离线脚本相比，Web 后端的 STFT 实现在矩形窗之外增加了 Hann、Hamming、Blackman、Bartlett 和 Kaiser 五种窗函数，并支持独立设置 FFT 点数（零填充），为用户提供了更灵活的参数空间。动画渲染通过异步任务队列完成，避免长时间阻塞 HTTP 请求，前端可轮询任务状态后下载生成的 MP4 文件。

\section{结论}
本文围绕短时傅里叶变换的可视化展开了完整实验。首先，通过 \texttt{waveform.py} 和 \texttt{fft\_analyze.py} 对两路给定信号进行时域观察与全局傅里叶变换，发现二者频谱几乎一致，说明普通频谱只能描述"频率集合"，无法反映"频率出现顺序"；其次，通过 \texttt{stft\_analyzer.py} 实现的矩形窗 STFT，对两路信号在不同窗长下的时频图进行了比较，验证了时间分辨率与频率分辨率之间的权衡关系，明确展示了两路信号在局部时序结构上的差异；再次，对自选音乐分别生成了滚动时频动画（\texttt{stft\_animator.py}）和彩虹雨频谱动画（\texttt{fft\_animator.py}），从不同角度展示了音频信号的动态频率结构；最后，将全部分析功能封装为 Flask 后端 API，结合 Vue.js 前端构建可交互网站并在线部署，使任意音频文件的时频分析可在浏览器中直接完成。

综上，STFT 是连接时域与频域的重要工具。对于非平稳信号，仅用傅里叶变换进行全局分析是不充分的；将滑动窗机制引入频谱计算后，才能获得更符合感知与工程需求的时间-频率联合表示。

\appendix
\section{核心代码}

\subsection{手写数值傅里叶变换（\texttt{lib/manual\_fft.py}）}

\begin{lstlisting}[language=Python]
def fft(signal, time_samples, frequency_samples):
    """按连续傅里叶变换的数值求和公式计算频谱。"""
    values = np.asarray(signal, dtype=np.complex128).reshape(-1)
    time_axis = np.asarray(time_samples, dtype=np.float64).reshape(-1)
    freq_axis = np.asarray(frequency_samples, dtype=np.float64).reshape(-1)
    sample_interval = float(time_axis[1] - time_axis[0])
    result = np.zeros(freq_axis.size, dtype=np.complex128)
    for k in range(freq_axis.size):
        for i in range(time_axis.size):
            result[k] += (
                values[i]
                * np.exp(-2j * np.pi * freq_axis[k] * time_axis[i])
                * sample_interval
            )
    return result
\end{lstlisting}

\subsection{STFT 核心实现（\texttt{stft\_analyzer.py}）}

\begin{lstlisting}[language=Python]
def stft(signal, fs, window_size, frame_shift=None):
    """矩形窗 STFT，默认半窗重叠，逐帧调用 numpy FFT。"""
    samples = np.asarray(signal, dtype=np.float64).reshape(-1)
    if frame_shift is None:
        frame_shift = max(window_size // 2, 1)
    nfft = window_size
    num_frames = (samples.size - window_size) // frame_shift + 1
    window = np.ones(window_size, dtype=np.float64)  # 矩形窗
    spectrum = np.zeros((nfft, num_frames), dtype=np.complex128)
    frequencies = np.arange(nfft, dtype=np.float64) * fs / nfft
    for m in range(num_frames):
        start = m * frame_shift
        frame = samples[start:start + window_size] * window
        spectrum[:, m] = np.fft.fft(frame, n=nfft)
    times = (np.arange(num_frames) * frame_shift + window_size / 2) / fs
    return frequencies, times, spectrum
\end{lstlisting}

\subsection{彩虹雨效果（\texttt{fft\_animator.py}）}

\begin{lstlisting}[language=Python]
def _build_rainbow_rain_segments(rng, min_frequency, max_frequency,
                                  min_db, max_db):
    """每帧生成 120 条对数均匀分布的 HSV 彩色线段作为装饰背景。"""
    drop_count = Config.rain_drop_count  # 120
    log_x = rng.uniform(
        np.log10(min_frequency), np.log10(max_frequency), drop_count
    )
    x_positions = np.power(10.0, log_x)
    y_starts = rng.uniform(min_db, max_db, drop_count)
    lengths = rng.uniform(
        Config.rain_length_min, Config.rain_length_max, drop_count
    )
    y_ends = np.maximum(y_starts - lengths, min_db)
    colors = plt.cm.hsv(rng.random(drop_count))
    widths = rng.uniform(
        Config.rain_width_min, Config.rain_width_max, drop_count
    )
    segments = np.stack((
        np.column_stack((x_positions, y_starts)),
        np.column_stack((x_positions, y_ends)),
    ), axis=1)
    return segments, colors, widths
\end{lstlisting}

\end{document}
